\begin{helper}[Assembly of a retained good hierarchy]
Let $G$ be a finite loopless multigraph with palette $[k]$,
distinguished edge $e_0$, and exact edge identities.
Fix an ETT $\mathcal E$ on $T$ generated by $\varphi$, a
\noderef{HierarchyBase} $B$, an arbitrary
\noderef{TashkinovHierarchy} $H$, and a witness
$g:\noderef{GoodHierarchy}(H)$.  Write $n=\mathcal E.\mathit{rounds}$,
$T_n,D_n$ for its terminal ladder tree and computed terminal stored set,
$q=H.\mathit{dividerCount}>0$, and
$A=\noderef{HierarchyStarredStage}(H,q)=H.\mathit{stage}(q)$.
For $0\le j\le q$ put $D_{n,j}=\noderef{HierarchyStoredColorsAt}(H,j)$,
and put
\[
 R_0=\bigcup_{\eta\in D_{n,0}}g.\Gamma^0_\eta.
\]
Let $\alpha\ne\beta$, let $z\in V\sqcup E$, and let
$\sigma=\noderef{TwoColorComponentSwap}(G,\varphi,\alpha,\beta,z)$.
Assume both colorings proper and leaving exactly $e_0$ uncolored.

Suppose supplied an input coloring, an actual certified ETT
$\mathcal E_A$ on the exact tree $A$ generated by $\sigma$, a base
$B_A:\noderef{HierarchyBase}(\mathcal E_A,\sigma)$, and an ordinary
hierarchy $K:\noderef{TashkinovHierarchy}(\mathcal E_A,B_A)$,
not assumed good.  Assume:
\begin{enumerate}
\item The terminal ladder tree, terminal extension kind, and computed
terminal stored set of $\mathcal E_A$ equal respectively $T_n$, the
terminal kind of $\mathcal E$, and $D_n$.
\item $K.\mathit{dividerCount}+1=q$, $B_A.\mathit{core}=B.\mathit{core}$,
$B_A.\mathit{closure}=B.\mathit{closure}$, and
$K.\mathit{stage}(j)=H.\mathit{stage}(j)$ for every $0\le j\le q$.
\item Every displayed edge of $A$ has equal values under $\sigma,\varphi$,
and every vertex of $A$ has equal missing-color sets under these colorings.
\item For every $0\le r\le q$ and every $c\in\{\alpha,\beta\}$ missing
and singleton-closed on $\noderef{HierarchyStarredStage}(H,r)$ under
$\varphi$, the recoloring is an \noderef{OutsideKempeRecoloring} of that stage.
\item If the original terminal kind is PE, there are a color
$c\in\overline\varphi(B.\mathit{core})\setminus R_0$ and an exact edge
$f\in\noderef{BoundaryColorEdges}(G,\varphi,V(B.\mathit{core}),c)$
with $\sigma(f)=\varphi(f)$.
\end{enumerate}
Then there is a witness $g_A:\noderef{GoodHierarchy}(K)$ such that
for every $0\le j<q$ its stored set is exactly $D_{n,j}$ and
\[
             g_A.\Gamma^j_\eta=g.\Gamma^j_\eta
                    \quad\text{for every }\eta\in[k],
\]
including colors outside the stored set.  Furthermore,
\[
 \noderef{TreeColorSetMinusClosed}
 \left(G,\sigma,H.\mathit{stage}(q),
       \bigcup_{\eta\in D_{n,q}}g.\Gamma^{q-1}_\eta\right)
\]
holds with the original last-stage stored set and preceding Gamma family.
\end{helper}
\begin{proof}
The exact nesting equations of \noderef{TashkinovHierarchy} show
$H.\mathit{stage}(j)\preccurlyeq A\preccurlyeq T$ for $0\le j\le q$:
compose consecutive prefix relations and use the terminal-stage equation.
The core lies in the first stage, hence in $A$, and in every starred
stage.  By \noderef{HierarchyBase}, $T_n$ is a prefix of the core
in non-PE by equality and in PE by the left-prefix clause of
\noderef{OrderedTreeUnion}.  In PE its vertex-union clause also puts
the closure's vertex set in the core.  Thus $T_n$, the core, every retained
ordinary and starred stage, and, in PE, the closure have vertex sets
contained in $V(A)$; the exact-prefix inclusions induce vertex-set
inclusions by \noderef{TreeSequenceInitialSegmentVertexSetSubset}.

For any exact prefix $U\preccurlyeq A$, take the union of the assumed
equal vertex missing sets to obtain
$\overline\sigma(U)=\overline\varphi(U)$ by \noderef{TreeMissingColors}.
The same union argument works for the PE closure using its vertex-set
inclusion.  Core equality and the stage equations identify every starred
stage of $K$ at $j<q$ with the old starred stage at $j$, by
\noderef{HierarchyStarredStage}.  The assumed equality of computed
terminal stored sets and the preceding missing-set equalities give
$D^K_{n,j}=D_{n,j}$ after expanding \noderef{HierarchyStoredColorsAt}.
The terminal ladder-tree equality gives equality of the zero-stage
missing sets, and the other stage equations give equality of all positive
stage-missing sets in \noderef{HierarchyStageMissingColors}.

Define the entire new Gamma family, not only its values at stored colors,
by $g_A.\Gamma^j_\eta=g.\Gamma^j_\eta$ for $j<q$.
We verify all law fields of \noderef{GoodHierarchy} for this family.
For two-color cardinality, a stored color at index $j$ belongs to exactly
the old stored set; apply the old two-color field.
For disjoint pairs, two distinct new stored colors are the same old
stored colors, so the old pairwise-disjointness field applies.

For avoidance at retained index $j<q$, put
$U=H.\mathit{stage}(j+1)=K.\mathit{stage}(j+1)$ and
$L=|\noderef{HierarchyStarredStage}(H,j)|$.
These are exact prefixes $U\preccurlyeq A\preccurlyeq T$.
Apply \noderef{P2813_T1_PrefixSegmentTransport} with the stated
edge-value and vertex-missing-set equalities, color $\eta$, and lower
cutoff $L$.  It identifies the new segment window, computed using the
totalized first-missing position on $A$ under $\sigma$, with the old
window computed on $T$ under $\varphi$.
Expanding \noderef{HierarchySegmentColors} therefore identifies the
two segment color sets.  This use includes the last-position fallback
when $\eta$ has no missing witness on $A$; equality of full first-missing
positions on $A$ and $T$ is not asserted.
The old avoidance field now puts each copied Gamma color in the unchanged
stage-missing set and outside the new segment color set.

For unique refresh at any internal index of $K$, the exact stage
difference equals its old counterpart.  Each vertex of that difference
belongs to $A$, so taking the union of its unchanged vertex missing sets
identifies the new difference-missing set with the old one.
The stored set and both compared total Gamma families agree as well.
The predicate selecting the refreshed stored color is consequently
identical to the old predicate, including its equality requirements for
every other stored color.  The old existence and uniqueness prove the
new field.

Consider the PE base field.  Equality of the actual terminal kinds makes
it vacuous if the original kind is not PE.  In PE choose the supplied
$c,f$.  The color remains missing on the core and outside $R_0$.
By \noderef{BoundaryColorEdges}, boundary membership of $f$ depends only
on endpoints, while its new color is still $c$ by the supplied value
equality.  New minus-closedness would force the $c$-boundary empty by
\noderef{TreeColorSetMinusClosed} and \noderef{TreeColorSetClosed},
contradicting this exact edge.  Stored-set and Gamma equality identify
the new initial reserve with $R_0$, so this proves the required
non-minus-closedness field.  For its second conjunct, $T_n$ and the PE
closure have unchanged missing sets, and the initial reserve is unchanged.
Thus the cardinality of their common missing colors outside the reserve
equals the old cardinality, which is at least four by the old PE base field.

Finally fix any $1\le j\le q$.  Old good-hierarchy minus-closedness applies
to $H.\mathit{stage}(j)$ with reserve
$R_j=\bigcup_{\eta\in D_{n,j}}g.\Gamma^{j-1}_\eta$.
Its vertex missing sets are unchanged by inclusion in $A$.
The stage is exactly the starred stage at $j$, so the assumed
closed-stage outside control applies to every affected color missing
and singleton-closed there.  These, the swap equality, properness,
one-uncoloredness, and $\alpha\ne\beta$ are precisely the premises of
\noderef{P2813_T1_MinusClosedTransport}.  It proves minus-closedness
under $\sigma$ with the same $R_j$.
For $j<q$, stored-set and copied-Gamma equalities make these assertions
all the internal minus-closed fields of $K$.
For $j=q$ the assertion is the separate terminal conclusion, using
the original $D_{n,q}$ and $\Gamma^{q-1}$ rather than the terminal
stored-set convention of $K$.  This constructs every good-hierarchy
field and proves both stated conclusions.
\end{proof}
